% Generated draft: requires mathematical and source review.
\documentclass[11pt]{article}
\usepackage[margin=25mm]{geometry}
\usepackage{fontspec}
\setmainfont{Latin Modern Roman}
\usepackage{amsmath,amssymb,mathtools,unicode-math}
\setmathfont{Latin Modern Math}
\usepackage{enumitem,needspace,xurl,hyperref}
\setlength{\parindent}{0pt}
\setlength{\parskip}{5pt}
\setlength{\emergencystretch}{3em}
\newcommand{\N}{\mathbb{N}}
\newcommand{\Z}{\mathbb{Z}}
\newcommand{\Q}{\mathbb{Q}}
\newcommand{\R}{\mathbb{R}}
\newcommand{\C}{\mathbb{C}}
\newcommand{\F}{\mathbb{F}}
\newcommand{\A}{\mathbb{A}}
\newcommand{\PP}{\mathbb{P}}
\newcommand{\E}{\mathbb{E}}
\newcommand{\eps}{\varepsilon}
\newcommand{\dd}{\,\mathrm{d}}
\begin{document}
\section*{UnsolvedMath ID: 30001076}
\textbf{Source problem code:} catalogue entry
\subsection{Definitions and mathematical statement}
Let $G$ be the particular complex-hyperbolic reflection group in the source's construction on the complex $13$-ball $\mathbb B^{13}$, let $\mathcal H$ be its mirror arrangement, and put $Y=\mathbb B^{13}\setminus\mathcal H$ and $X_0=Y/G$. The extract does not give the lattice or generators specifying $G$, so this is not a statement about an arbitrary ball reflection group. Let $\Pi$ denote the fundamental group intended by the source's notation $\pi_1(X_0)$. Ordinary versus orbifold fundamental group must be specified: residual stabilizers of the action can distinguish them. For the indicated meridian loops $m$ about discriminant components, the intended order-two quotient is schematically \[Q=\Pi/\langle\!\langle m^2\text{ for the indicated meridians }m\rangle\!\rangle.\] This is a quotient of the appropriate discriminant-complement group, not a silent simultaneous imposition of $m^2=1$ and an already-existing $m^3=1$ relation, which would kill $m$. The source's phrase order two rather than order three needs the precise construction to determine the group and relations being replaced.\par Let $\mathbb M$ denote the Monster simple group. The bimonster is $(\mathbb M\times\mathbb M)\rtimes C_2$, where the nonidentity element of $C_2$ exchanges the two factors. The first question is whether \[Q\simeq(\mathbb M\times\mathbb M)\rtimes C_2.\] The report says the proposed quotient satisfies its known bimonster relations except possibly $w^{10}=1$, but does not provide the word $w$ or the relation list in this extract. The additional questions are whether the source's sixteen corresponding based loops generate $\Pi$, and what further relations $\Pi$ has. Generation of the quotient $Q$ is not substituted for generation of $\Pi$ itself. The sixteen loops, their base paths, and their configuration are not specified. The supplied 2026-08-21 triage labels the record partially solved and reports that a 2025 primary paper reduces the quotient identification to two possibilities, the bimonster or a group of order two, without settling which. That progress report is retained without treating the final identification or the separate sixteen-loop presentation questions as resolved.
\subsection{Short English statement}
Does the order-two meridian quotient of the source's ball-arrangement group equal the bimonster, and do its sixteen specified loops generate the full complement group? The reflection group, meridians, fundamental-group convention, and exceptional word are missing from the extract.
\subsection{Sources}
\begin{itemize}\item Local UnsolvedMath registry, permanent ID 30001076, source problem catalogue entry, \emph{Bimonster Quotients of Complex Hyperbolic Reflection Groups}, proposed year 2008; no individual proposer supplied.\item Supplied citation: \emph{Geometry and Arithmetic around Hypergeometric Functions} (2009), Oberwolfach Reports, DOI 10.4171/owr/2008/45, pp.\ 2557--2620: \url{https://doi.org/10.4171/owr/2008/45}. Supplied literature label: OWR Report 45.\item Supplied literature label: The Deligne--Mostow 9-ball and the monster, \url{https://doi.org/10.2140/gt.2025.29.791}.\item Supplied literature label: Allcock--Basak preprint, \url{https://arxiv.org/abs/2309.00148}.\item Literature assessment date: 2026-08-21. The status and progress claims, including the reported 2025 reduction to two possibilities, are supplied dated curation, not independently verified here.\end{itemize}
\end{document}
